\section{吴越：江水尽头的两种雄心}

从淮河向南，车道常在水边断开，军队须把粮秣和兵器装上船；从太湖到钱塘江，港汊、湖沼与支流又把一座城同另一座城隔开。这不是中原会盟的边角，而是一套能养舟师、容纳聚落、也能遮蔽行军消息的地理。吴与越进入《春秋》《左传》的视野时，北方诸侯已经习惯以车战、盟书和周王室名分谈论天下；下游的两国却能借水陆相接的通道，把自己的兵力送到江汉与中原。它们并非只是接受那套语言，也改变了谁能在其中发言。

水网也不是一份可以随时兑现的优势。船只使兵员、粮秣和消息能越过车道的断裂处，却要有港汊、城邑与能够征发的人力来接续；离开本土越远，补给和指挥就越依赖这些节点。寿梦即位后，吴开始以诸侯身份稳定进入中原材料的视野。\eventanchor{SA-0585-WU-SHOUMENG-ASCEND}{春秋·寿梦即位}传统系年的早期吴王世系仍须在《左传》与《史记》之间互校，但晋楚竞争为江淮下游的政权提供了被看见、也被动员的机会。前584年，\eventanchor{SA-0584-WU-JIN-CONTACT}{春秋·晋人通吴}把这种机会变成了关系：晋要在楚的东侧另开压力，吴也需要北方联系扩大行动范围。《左传》所说“教以车战”，应留在叙事传统的限度内，不能由此逐项复原一场军事革命；较可把握的是，联络使吴的水陆通道开始嵌入原先由晋、楚主导的竞争。

此后的若干节点显示，这不是一次使节往来便能完成的转向。楚军曾深入鸠兹至衡山，吴又伏击楚军并反取驾；经年大体可靠，具体行军和战斗细节则仍有分歧。\eventanchor{SA-0570-CHU-WU-JIUZI}{春秋·楚伐吴与鸠兹之战}前563年晋与吴会于柤、前559年士匄又会诸侯与吴于向，\eventanchor{SA-0563-ZHA-WU-MEETING}{春秋·晋吴会柤}\eventanchor{SA-0559-XIANG-MEETING}{春秋·晋吴会向}使吴得到在盟会中出现的席位，却也把它放进反楚网络的期待之下。前556年吴军沿江淮方向伐楚至豫章而退，地望与路线尚有不同解释；前537年楚联合蔡、陈、许、顿、沈、徐、越等再伐吴，越人身份和规模也不宜按后世吴越故事夸大。\eventanchor{SA-0556-WU-CHU-YUZHANG}{春秋·吴伐楚入豫章}\eventanchor{SA-0537-CHU-WU-COALITION}{春秋·楚再伐吴}这些来回说明，航道既让吴能试探楚的东部防线，也让楚必须把属国、淮夷和水陆节点编进防务。

吴没有只停留在被动防守。前519年鸡父之战击破楚方联军，前518年灭巢，前512年又灭徐；兵力、路线和控制程度多只能从《左传》叙事谨慎推定，但江淮间的城邑与航道节点确实逐步改换了归属。\eventanchor{SA-0519-WU-JIFU}{春秋·鸡父之战}\eventanchor{SA-0518-WU-CHAO}{春秋·吴灭巢}\eventanchor{SA-0512-WU-XU}{春秋·吴灭徐}前529年吴趁楚国内乱取州来，也使北上通道更有支点。\eventanchor{SA-0529-WU-ZHOULAI}{春秋·吴灭州来}这不是一条直线扩张：楚对徐的控制与进攻仍在持续，徐的处境正显示小国在楚、吴之间所承受的压力。\eventanchor{SA-0538-CHU-DETENT-XU}{春秋·楚执徐子}\eventanchor{SA-0530-CHU-ATTACK-XU}{春秋·楚伐徐}不过，吴已不再只是晋人用以牵制楚的一枚远方棋子。这个使能条件，终于在\eventref{SA-0506-BOJU}{春秋·柏举之战}中显出锋芒。

阖闾即位和伍子胥入吴，是吴向楚扩张的又一组织条件；伍子胥在早期材料中的位置较清楚，孙武与他分别承担何种谋划则多出于后出的传记安排。\eventanchor{SA-0514-WU-ZIXU}{春秋·阖闾即位与伍子胥入吴}前515年，吴利用楚国内部危机筹划北上、西进，长期整军而非一次偶发远征，构成了柏举的近因。\eventanchor{SA-0515-WU-HELÜ-CAMPAIGN}{春秋·阖闾筹划伐楚}

\begin{event}{公元前506年}{柏举之战，吴军攻入郢都}{可靠纪年}{柏举至郢都}\eventref{SA-0506-BOJU}{春秋·柏举之战}
楚军在汉水一带列阵，自小别、大别而至柏举；吴与蔡、唐联军连战推进。战事最初并不像一条笔直的进军线：楚令尹囊瓦见势不利而离军出奔，楚昭王随后离开郢都，吴军才得以沿楚军后路追击，进入这座江汉中心。郢都失守，是楚国罕有的深创，也是吴国第一次把下游的军力直接推到大国都城门前。
\end{event}

一座都城的陷落很容易使胜者相信，整个国家已经在掌中。可是攻进郢都与支配江汉是两回事。吴军的补给线拖在远处，楚的腹地、旧有的地方关系与仍可调动的人力却没有随着城门一同消失。楚昭王的出奔、秦的援助以及吴军难以长期停留，共同限制了胜利的转化；次年楚政权已开始收复核心地区。柏举并未使吴无足轻重，恰恰相反，它使各国不得不把吴算作天下竞争者；但它也预先显露一个新兴国家的难题：远征能够击穿对手，未必能够把陌生而辽阔的土地编入自己的征发与治理。

\begin{textwindow}[柏举、伍子胥与孙子的名字]
《左传·定公四年》以楚军部署、囊瓦出奔、昭王出走和吴入郢构成柏举的叙事骨架，文本中没有孙武。伍子胥在吴楚冲突中的位置则更为清楚：他由楚至吴，其家族遭遇与对楚用兵相连，后来的《史记》尤其把这条线写成复仇与功业交织的故事。司马迁又在《孙子吴起列传》中赋予孙武显著位置；传世《孙子兵法》的早期核心及银雀山汉简，证明这部兵书有很早的文本传统，却不能反证孙武必在柏举现场统军。将孙子当作后世解释吴军胜利的重要传统可以；把他写成由最早材料确证的战场指挥者则不可以。
\end{textwindow}

吴的西进尚未结束，南方的难题已经逼近。前510年吴伐越，吴越直接冲突才明确进入《春秋》经文；战役规模和双方意图不足以让人倒推后来那种全面争霸。\eventanchor{SA-0510-WU-YUE}{春秋·吴伐越}前496年槜李之战中，越败吴，阖闾伤重而死；经年可靠，战场细节却有争议。\eventanchor{SA-0496-ZUILI}{春秋·槜李之战}夫差继位后首先要回答的，正是这场未竟战争。\eventanchor{SA-0496-FUCHAI}{春秋·夫差即位}对吴而言，越据有南岸水网，不能被当作身后的小患；对越而言，吴既握有更强的兵力，又能调动北方声望，正面相持的余地极小。前494年，夫差在夫椒击败勾践。\eventanchor{SA-0494-FUJIAO}{春秋·夫椒之战}夫椒的确切地点与战场展开难以从传世材料中精密钉定，但败后越国以会稽为最后依托、遣大夫向吴求成的基本线索，见于《左传·哀公元年》。

《左传》使伍子胥反对受降，太宰嚭则主张纳越，这一组相反意见把吴廷的选择写得格外鲜明；《国语·吴语》《越语》与《史记·越王勾践世家》又为它添上更多谋臣、进献和屈辱的场景。可以把它们读作失败者如何获得喘息的政治叙事，却不必把每一句劝说都当成原始会议纪录。吴接受和约，短期看是以较低成本消除南方威胁，得以转用兵力；它付出的代价是让越国保存了君主、贵族网络和重新组织资源的时间。勾践是否一度亲至吴廷服役、范蠡和文种各自究竟提出了哪些方案，材料的铺陈程度并不相同；较可把握的是，越没有在夫椒之后消失。

归国后的重建，也不应被压缩成一个人苦熬的传奇。越要修复的首先是服从关系：败战后还能征集谁，谁能守住水道和城邑，怎样让贵族、官吏与聚落相信国家仍有可供依附的中心。农业生产、兵员训练、馈赠往来和对吴的臣服姿态，都是重建空间的一部分；《国语·越语》所说“十年生聚，十年教训”把这个长过程凝结成整齐的战略格言，适合提示恢复并非一朝一夕，却不能替代逐年可核的行政档案。前490年代初，越已逐渐恢复对外行动的能力，\eventanchor{SA-0490-YUE-RECOVERY}{春秋·越国重建}所指的正是这一有明确方向、却难以精确切成年表条目的过程。

吴廷并非无人看见这一变化。越在重建，吴却继续把舟师、人力与外交注意力投入北方。前489年吴伐陈，前488年鲁哀公会吴于鄫，前487年吴又伐鲁而齐乘机取讙、阐二邑；经文所记的行动可靠，吴齐是否协调、吴军意图为何，不能写成确定的幕后安排。\eventanchor{SA-0489-WU-ATTACK-CHEN}{春秋·吴伐陈}\eventanchor{SA-0488-WU-ZENG-MEETING}{春秋·鲁吴会鄫}\eventanchor{SA-0487-WU-LU}{春秋·吴伐鲁}前485年鲁哀公会吴伐齐，吴军由鲁道北上；鲁为何参战、补给线怎样安排与兵力多少，材料不足以精确回答。\eventanchor{SA-0485-WU-QI-CAMPAIGN}{春秋·吴鲁伐齐}次年艾陵之战中吴鲁联军击败齐军、获齐国书，吴的北上威望到达高处；战术与俘获经过主要保存在《左传》叙事，不能把它们当作无缝的战场记录。\eventanchor{SA-0484-AILING}{春秋·艾陵之战}\eventanchor{SA-0484-WU-CAPTURES-GUOSHU}{春秋·吴获齐国书}

伍子胥的地位本来就建立在对楚战争、对外谋划和君主信任之上；当夫差把目光投向北方，他的经验使他更倾向于把越视为不能搁置的危险。可是一个功臣的警告能有多大分量，取决于君主、近臣与军功集团如何分配信息和赏赐。艾陵之后北上的道路似乎更加开阔；同年，伍子胥被夫差赐剑自尽。\eventanchor{SA-0484-WUZIXU-DEATH}{春秋·伍子胥之死}

《左传·哀公十一年》《国语·吴语》与《史记》都保存了君臣决裂的主线，却以不同的言辞、预言和临终姿态解释它。后世最熟悉的伍子胥，是掘墓鞭尸、悬目望越的复仇者；那些强烈的画面使人物不朽，也最需要同较早编年叙事分开。能够从多种材料的交集里提出的判断较为有限：吴失去了一位主张防越、又有对外战争资历的重臣，夫差则更能把北上所获的声名当作现实的政治资本。这不是说一个人的死注定了吴亡；而是说，朝廷里能够把短期胜利同南方风险相联系的声音少了一种制度位置。

艾陵后的北方关系还须维系：前483年鲁哀公会吴于橐皋，是吴继续经营鲁与东部道路的一个节点。\eventanchor{SA-0483-WU-TAOGAO}{春秋·鲁吴会橐皋}两年后，夫差终于把吴军与吴的名分带到黄池。黄池远在北方，会盟的礼序、歃血的先后与谁居盟主，都是可以被诸侯看见的权力。\eventanchor{SA-0482-HUANGCHI}{春秋·黄池之会}《左传·哀公十三年》所写的吴晋争先，不只是两位君主的意气相争：吴若在旧有的春秋礼法中居先，便仿佛证明下游国家已经取得组织诸侯的资格。夫差所追逐的并非空名；北方盟友、贸易通道、军事威望和对国内功臣的赏赐，都可能由这一席位增益。

然而会盟不能同时守住都城。夫差停留黄池之际，越军已向吴境发动进攻，吴的南方防线和王室承受了北上难以照看的压力。战事的先后、损失和使者言辞，在《左传》《国语》与《史记》中各有安排，不宜把它们拼成一份毫无缝隙的战报；它们共同指出的结构却很清楚：吴把有限的军政注意力分给了两个方向，越则不必在黄池的礼仪语言中取胜，只需在本土水陆之间持续积累优势。

\begin{causalchain}[从郢都到姑苏]
\eventref{SA-0506-BOJU}{春秋·柏举之战}使吴看见介入天下的道路；\eventref{SA-0494-FUJIAO}{春秋·夫椒之战}暂时给了它处置南方的机会。吴选择北上，并在\eventref{SA-0482-HUANGCHI}{春秋·黄池之会}取得可见的声望；越则把和约留下的存续空间转为重建与反攻的条件。这里没有一条由胜至败的宿命线，而是远征成本、朝廷判断和两国地理资源不断相互反馈，终于通向\eventref{SA-0473-WU-FALL}{春秋·越灭吴}。
\end{causalchain}

黄池之后，吴虽曾应战、求和，也没有立即覆亡；越的攻势同样需要粮食、舟师、盟友与持续动员。约前478年，越再攻吴，夫差在败势中求和；约前476年，越继续进攻，压缩吴都周边的回旋空间。\eventanchor{SA-0478-YUE-ATTACK-WU}{春秋·越再攻吴}\eventanchor{SA-0476-YUE-SECOND-ATTACK-WU}{春秋·越第二次攻吴}这两个年次来自传统年表，围攻过程与具体日期在《国语》《史记》及《竹书纪年》之间并不整齐，不能补成逐日战报。它们仍共同表明，吴的盟友、供给和防线已在持续进攻下承受压力。数年反复之后，吴在前473年为越所灭，夫差自尽，吴的王室政权退出了这场竞争。\eventanchor{SA-0473-WU-FALL}{春秋·越灭吴}\eventanchor{SA-0473-FUCHAI-DEATH}{春秋·夫差自尽}《史记·越王勾践世家》采用前473年，传世纪年与现代历法重建之间仍有前后数年的讨论，因此这一年份应视为通行的编年节点，而不是能精确到某日的结案。越一度北上会盟，勾践也成为诸侯政治中不能忽略的人物；但这并不意味着它继承了一个稳定的“南方霸权”。吴留下的水陆网络、北方介入的经验和战争损耗，都仍在改变下一个时代的权力地图。

\begin{turningpoint}[制度得失：声望的半径与防务的半径]
吴把水网、舟师和江淮通道转成远征能力，先后击楚、制越、伐齐、会黄池，解决了新兴国家难以被天下承认的问题。成本也随之累积：远征需要把兵员、粮秣和君主注意力带离本土，而会盟的名分不能替代南岸的驻守与国内的风险判断。越的胜利不是“忍辱”自动兑现的报偿；它来自败后仍能保存组织、在较长时间内重建资源，并抓住吴将安全与声望分置两端的时刻。
\end{turningpoint}

\begin{sourcenote}[“卧薪尝胆”与吴越故事的层累]
吴越故事尤其容易成为道德寓言。勾践“尝胆”的说法可见于《史记》的越王叙事，后来固定为“卧薪尝胆”的成语；“卧薪”的具体情节和两字连用则属于更晚的文本层累，不能当作春秋现场已有的政治仪式。《左传》最接近相关年代，却也以战争和外交为中心；《国语》的吴语、越语擅长把复杂的谈判塑成可传诵的谋略言辞；《史记》再将人物命运编成完整的兴亡剧。出土的越王兵器能够提示下游精英文化与铸造技术的高度，却不能替这些叙事逐一证明谋臣密计、君主苦行或战场细节。把不同材料放在各自的位置上，吴越的兴亡反而更有重量：它不是一句格言的奖惩，而是两个国家在河流、兵役、朝廷与时间中作出的选择。
\end{sourcenote}

\phantomsection\label{fig:vol01:spring-autumn:boju}
\begin{center}
\begin{tikzpicture}[x=.88cm,y=.88cm,font=\small]
  \fill[paper] (0,0) rectangle (10.5,5.7);
  \draw[blue!50,line width=.9pt] (.4,4.8) .. controls (2.5,4.6) and (4.3,4.9) .. (6.3,4.4);
  \node[text=blue!70!black,above] at (3.5,4.75) {淮河通道};
  \draw[blue!55,line width=1pt] (5.2,.4) .. controls (6.1,1.8) and (7.0,2.2) .. (9.8,1.5);
  \node[text=blue!70!black,below] at (7.3,1.75) {汉水与江汉平原};

  \fill[cyan!22] (.8,3.2) -- (3.4,3.0) -- (3.8,4.8) -- (1.2,5.1) -- cycle;
  \node[align=center] at (2.1,4.0) {吴\\下游水网};
  \fill[purple!20] (5.4,1.1) -- (9.8,1.0) -- (10.0,4.4) -- (6.3,4.6) -- cycle;
  \node[align=center] at (7.8,3.0) {楚\\江汉腹地};
  \fill[orange!20] (8.7,4.4) -- (10.1,4.4) -- (10.2,5.4) -- (8.9,5.5) -- cycle;
  \node at (9.5,4.95) {郢都方向};
  \fill[timeline] (5.2,3.5) circle (.14);
  \node[above,align=center] at (5.2,3.6) {柏举\\战场};

  \draw[-{Stealth[length=2mm]},ultra thick,color=cyan!60!black] (2.7,3.8) -- (5.05,3.5);
  \draw[-{Stealth[length=2mm]},ultra thick,color=cyan!60!black] (5.35,3.45) -- (8.5,4.75);
  \node[text=cyan!60!black,above] at (4.0,3.65) {吴军经北线深入};
  \node[text=cyan!60!black,left] at (7.1,4.15) {向郢推进};
\end{tikzpicture}

\smallskip
\small\textit{图  柏举战役空间示意：强调吴军利用北方通道进入楚地、柏举战场与郢都方向的关系；不表示精确战线、河道或疆界。}
\end{center}

